\documentclass[11pt]{article}

\usepackage[utf8]{inputenc}
\usepackage[T1]{fontenc}
\usepackage[margin=1in]{geometry}
\usepackage{hyperref}
\usepackage{amsmath,amssymb}
\usepackage{graphicx}
\usepackage{booktabs}
\usepackage[expansion=false]{microtype}
\usepackage[numbers]{natbib}
\usepackage{placeins}

\hypersetup{
    colorlinks=true,
    linkcolor=black,
    citecolor=blue,
    urlcolor=blue,
}

\title{\textbf{Code Correctness Is Linearly Decodable from LLM Hidden States Before Generation}}
\author{Carlo Di Cicco\thanks{Independent researcher. Contact: \texttt{cdicicco2001@gmail.com}. Code, data, and analysis scripts: \url{https://github.com/CarloDiCicco/ReasoningLab}}}
\date{}

\begin{document}

\maketitle

\begin{abstract}
Large language models encode rich information in their hidden states. 
This work asks whether the correctness of code that Qwen3-4B-Instruct-2507 has not yet generated is already legible in its hidden states, evaluated on a set of 444 tasks from LiveCodeBench.
The correctness of the model's first-attempt code is linearly decodable from the hidden state at the final prompt token, captured before any output token is generated, with a leakage-free held-out AUC of $0.881 \pm 0.008$ across 50 outer splits. 
To assess whether this signal is explained by prompt length, each hidden state dimension is residualized with respect to its linear effect. The probe still achieves an AUC of $0.842 \pm 0.010$, substantially above a logistic prompt-length baseline of $0.657 \pm 0.014$, and none of the nonlinear models tested improves upon it. 
A companion question about whether self-repair leaves a geometric signature in the model's hidden states could not be answered, because successful repairs following a failed first attempt are too rare in this setting to support the analysis. 
The contribution is both empirical and methodological, providing evidence that pre-generation hidden states contain a robust signal of eventual code correctness, together with a confound-control diagnostic that quantifies how much of that signal survives adjustment for prompt length.
\end{abstract}

\section{Introduction}
\label{sec:introduction}

Reasoning has become a central axis of large language model development. With each generation, models train on more and better data with improved training methods, and become capable of tackling problems that would have been out of reach a few years ago. Their internals, in turn, grow more informative: a richer internal computation is doing the work behind richer external behavior, which makes the question of what those internals encode increasingly worth asking. Code generation is a particularly clean field for that question: outputs are easily verifiable against unit tests, with no human-judgment loop required to decide whether an answer is correct. As model performance on code-generation benchmarks has improved rapidly, the question is no longer only \emph{whether} a model can produce correct code, but also \emph{what happens internally} as it does so. This shifts the focus toward mechanistic interpretability in models that now exhibit genuine capability in code generation.

A standard tool for asking what is encoded inside a frozen language model is the \emph{linear probe}: a simple linear classifier trained on a hidden state vector to predict some property of interest, where probe performance is read as a measure of how linearly decodable that property is from the representation \citep{alain2016probes}. Two recent works have studied code correctness through hidden states \citep{bui2025openia, ribeiro2025internal}, reading hidden states of already-generated code to judge whether it is correct without running the unit tests. One trains a classifier on those states; the other contrasts correct and incorrect samples to extract a correctness direction and rank candidate generations. In both, the hidden state is read \emph{after} the model has produced an answer.

This work asks a related but earlier question: whether correctness is already decodable from the hidden state at the position of the last prompt token, captured on a single forward pass over the prompt and before any output token is sampled. Whereas those works operate on hidden states of generated code, the probe in this paper operates on the hidden state of the prompt alone. This is a more upstream interpretability claim: the information must already be present at the moment the model finishes reading the input, not just at the moment it has committed to a particular output.

A second, more geometric question follows. When the model is given the chance to repair its own failing code, does the shift in its hidden state between the failed attempt and the repair attempt carry a stable signature of whether the repair will succeed? Phrased this way, the question lives within the Representation Engineering framework \citep{zou2023repe}, which identifies candidate features as \emph{directions} in activation space, typically as the difference between two group means. The natural application here is to compute the per-task hidden state delta from a failing attempt to its repair attempt, split the deltas by repair outcome, and ask whether the contrastive direction is stable.

Such contrastive constructions, however, are subtler than they look. A direction of the form $\mathrm{mean}(\text{success}) - \mathrm{mean}(\text{failure})$ already cancels any covariate whose mean is equal across the two groups, so residualizing the deltas against such a covariate would only collapse the signal artifactually. Conversely, when a covariate's group means differ, the raw direction is partially driven by it rather than by the latent feature of interest. The methodological approach of this paper makes this effect explicit: a short distribution test per candidate covariate, run before any residualization, decides which residualizations are warranted and which would be misleading.

The two findings are the following. The pre-generation probe holds up cleanly: the prompt-final hidden state carries a linearly decodable correctness signal that is robust across many random train/test splits and that survives controlling for prompt length. The repair direction is likewise statistically detectable in the raw deltas, passing both magnitude and inclination tests against label-shuffled nulls, and it remains significant after a conditional residualization once three observable repair-context covariates are checked and found to differ between the success and failure groups, even as that residualization removes a substantial part of it. The honest reading is that the geometric direction is a genuine repair-success signal that the observable repair context does not fully account for, rather than a mere artifact of that context. Consistent with prior reports that iterative self-repair plateaus after the first one or two attempts \citep{arimbur2026howmanytries}, the per-attempt pass rate collapses sharply after the first repair, and later transitions are too data-starved to support the same direction analysis.

Read together, the two results illustrate how the same confound control procedure applies to two different signals. For the pre-generation correctness probe, residualizing against prompt length confirms that the decodable signal is not reducible to prompt structure, even if a small part of the raw probe's performance is. For the geometric repair direction, residualizing against the differentially distributed repair-context covariates removes the part of the signal that the repair context can account for, and a repair-success signal remains once that part is taken out. In both cases the procedure measures how much of a raw signal is genuinely its own.
\section{Setup}
\label{sec:setup}

\subsection{Model and data}

All experiments use \texttt{Qwen/Qwen3-4B-Instruct-2507} \citep{qwen3report2025}, a 36-layer decoder-only transformer with a 2560-dimensional hidden state. The dataset consists of 444 LiveCodeBench \citep{jain2024livecodebench} tasks spanning the easy, medium, and hard difficulty tiers. Every experiment runs on a single NVIDIA DGX system.

\subsection{Hidden state capture}

Across all task attempts, the hidden state at the final prompt token is extracted from a prompt-only forward pass, before the autoregressive continuation begins. This yields a single 2560-dimensional vector per task per attempt at each captured layer. Hidden states are captured for the upper transformer blocks (layers 29--36), motivated by the general finding that linear separability of high-level properties increases with depth in deep networks \citep{alain2016probes}. Layer selection is handled by nested cross-validation (CV) in Section~\ref{sec:probe_results}.

\subsection{Repair-B policy}

For tasks where the model's first attempt fails the unit tests, an iterative repair procedure (``repair-B'') is invoked, with a maximum of 5 total attempts per task. Each attempt is a fresh forward pass on the full prompt, with no state carried over from previous attempts, and the repair prompt contains the original problem, the model's previous code, the error output captured by the verifier, and an instruction to fix the failure. Generation is capped at 1024 new tokens per attempt. The pre-generation probe in Section~\ref{sec:probe_results} uses all 444 tasks at attempt 0, and the outcome of the repair procedure itself is summarized in Section~\ref{sec:limitations}.

\section{Correctness is linearly readable before generation}
\label{sec:probe_results}
Before any token is generated, the hidden state at the last prompt token encodes information about whether the model's first-attempt generation will be correct. With the probe layer selected by nested CV, as outlined in Section~\ref{subsec:layer_selection}, a linear probe trained on this single vector predicts pass/fail on held-out tasks with mean test area under the ROC curve (AUC) of $0.931 \pm 0.008$ over $50$ random splits, well above a prompt-length-only baseline of $0.754 \pm 0.014$.

\subsection{Procedure}

For each of the 444 LiveCodeBench tasks, the hidden state at the last prompt token of the first attempt (attempt 0) is captured, producing a single 2560-dimensional vector per task per layer, as described in Section~\ref{sec:setup}. Each vector is labelled by whether the attempt-0 code passed the task's unit tests, and the probe is trained to predict this label. The pipeline is a three-step composition:
\begin{enumerate}
    \item \texttt{StandardScaler}, fit on the training fold only.
    \item \texttt{PCA} retaining 95\% of cumulative variance, fit on the training fold only.
    \item \texttt{LogisticRegression} with $\ell_2$ regularization. The regularization strength $C$ is selected by 5-fold CV inside the training fold, scoring on \texttt{roc\_auc}, over the grid
    \[
        \{10^{-3}, 10^{-2}, 10^{-1}, 1, 10\}.
    \]
\end{enumerate}

All three pipeline components are implemented in scikit-learn \citep{pedregosa2011sklearn}.

To get a stable estimate rather than rely on a single split, the pipeline is fit and evaluated $50$ times under different random stratified $80/20$ splits, and performance is summarized by the mean test AUC with a $95\%$ confidence interval. Three configurations are compared: two apply the three-step pipeline above to different feature sets, and the third is a simpler prompt-length baseline.
\begin{enumerate}
    \item \textbf{Raw}: the three-step pipeline applied directly to the hidden state vector as captured.
    \item \textbf{Residualized}: for each of the 2560 hidden state dimensions, an OLS regression is fit on the train fold with that dimension as response and \texttt{prompt\_tokens} as the single predictor. The fitted prediction is subtracted from both train and test for that dimension, and the standard three-step pipeline is applied to the residualized features. Residualization is therefore re-fit on every random split.
    \item \textbf{Prompt-only baseline}: a logistic regression on \texttt{prompt\_tokens} alone (a single feature, so no dimensionality reduction as PCA is needed), fit and evaluated under the same 50 splits with $C$ selected by cross-validation as above.
\end{enumerate}

The residualized configuration and the prompt-only baseline both target the same confound. In LiveCodeBench, prompt length is itself correlated with task difficulty: harder problems tend to come with longer specifications. A probe trained on raw hidden states could in principle exploit this by reading off prompt length rather than anything mechanistic about the model's internal representation. Residualizing each hidden state dimension against \texttt{prompt\_tokens} removes the linear contribution of prompt length from every dimension, leaving only the part of the hidden state that is not linearly explained by it. The prompt-only baseline goes one step further: it asks how well a probe can do using \texttt{prompt\_tokens} \emph{alone}, giving a reference the residualized probe must clear to show it carries signal beyond length.

The single sample per task ensures no task appears in both train and test splits, removing the inflation that would arise from leakage across attempts of the same problem.

\subsection{Layer landscape and selection}
\label{subsec:layer_selection}

The raw three-step pipeline is applied to each of the captured upper-block layers (29--36). Running it independently per layer, on 50 random stratified $80/20$ splits each, gives the descriptive landscape in Table~\ref{tab:layer_sweep}: a gentle plateau that peaks at the lower bound of the captured range (layers 29--30) and declines slightly toward layer 36. Unlike the monotonic increase of linear separability with depth reported for convolutional networks \citep{alain2016probes}, the probe AUC here peaks before the final block and then declines, suggesting that final transformer layers shift from maintaining general semantic representations to optimizing for next-token prediction.

\begin{table}[h]
\centering
\begin{tabular}{lcccccccc}
\toprule
Layer & 29 & 30 & 31 & 32 & 33 & 34 & 35 & 36 \\
\midrule
Test AUC & $0.934$ & $0.932$ & $0.924$ & $0.920$ & $0.920$ & $0.919$ & $0.917$ & $0.913$ \\
CV AUC   & $0.938$ & $0.939$ & $0.933$ & $0.929$ & $0.929$ & $0.926$ & $0.925$ & $0.923$ \\
\bottomrule
\end{tabular}
\caption{Per-layer probe performance on the 444 LiveCodeBench tasks: an independent per-layer sweep, using its own 50 random stratified 80/20 splits, each layer is evaluated on both mean test and CV AUC. The upper blocks form a shallow plateau peaking at layers 29--30. This table is descriptive only: selecting the argmax layer on these test scores would constitute selection on the test set, which the nested CV procedure below avoids.}
\label{tab:layer_sweep}
\end{table}

Because the layer is a hyperparameter, choosing it by the test-set argmax of Table~\ref{tab:layer_sweep} would leak the test set into the model-selection step and bias the reported performance upward \citep{cawley2010overfitting}. The layer is therefore selected by nested CV, following the same validation-based layer-selection logic used in prior work on internal correctness representations \citep{ribeiro2025internal}. Across 50 outer stratified $80/20$ splits (the outer loop, whose held-out folds provide the unbiased estimate), the layer, jointly with the regularization strength $C$, is selected by 5-fold CV within each outer training fold; the selected configuration is refit on the full outer-training fold and evaluated once on the untouched outer-test fold. Layer 29 emerges as the modal selection, chosen in $25$ of $50$ outer splits, with layer 30 a statistical near-tie at $24$; Table~\ref{tab:nested_cv_selection} shows that selection concentrates almost entirely on these two layers. The configuration selected most frequently across the 50 splits is layer 29 and $C = 0.01$, chosen in $22$ of them; this modal frequency is the criterion by which the final hyperparameter combination is fixed. The leakage-free outer-test AUC of the full select-then-train procedure is $0.931 \pm 0.008$. Layer 29 is adopted for all subsequent analyses on the basis of this modal selection; the descriptive sweep in Table~\ref{tab:layer_sweep} is consistent, ranking layer 29 highest ($0.934$), with layer 30 statistically indistinguishable.

\begin{table}[h]
\centering
\begin{tabular}{lcccccccc}
\toprule
Layer & 29 & 30 & 31 & 32 & 33 & 34 & 35 & 36 \\
\midrule
Times selected (of 50) & $25$ & $24$ & $1$ & $0$ & $0$ & $0$ & $0$ & $0$ \\
\bottomrule
\end{tabular}
\caption{Nested-CV layer selection: how often each layer is chosen, across all $C$ values, as part of the best (layer, $C$) configuration over the 50 outer splits. The grid search operates jointly over layer and $C$; this table shows the layer counts collapsed across $C$ to make the layer preference visible. Selection concentrates on layers 29--30, accounting for $49$ of $50$ splits, with deeper layers never selected; within these, $C = 0.01$ is the most frequent regularization strength, appearing in $43$ of $50$ splits. The full procedure's leakage-free outer-test AUC is $0.931 \pm 0.008$.}
\label{tab:nested_cv_selection}
\end{table}

\subsection{Results}

Each number in Table~\ref{tab:probe_results} is a leakage-free procedure-level estimate: for the raw and residualized probes, layer and $C$ are jointly selected inside each outer training fold and evaluated on the untouched outer-test fold; for the prompt-only baseline, only $C$ is selected as there is no layer. Reporting procedure-level numbers for all three configurations ensures they are directly comparable, since each reflects the performance of the full select-then-fit pipeline applied consistently, rather than locking any configuration to a specific hyperparameter value chosen by inspecting the test set.

\begin{table}[h]
\centering
\begin{tabular}{lc}
\toprule
Configuration & Test AUC (mean $\pm$ 95\% CI) \\
\midrule
Raw hidden states    & $0.931 \pm 0.008$ \\
Residualized         & $0.911 \pm 0.010$ \\
Prompt-only baseline & $0.754 \pm 0.014$ \\
\bottomrule
\end{tabular}
\caption{Nested CV procedure-level performance for the three probe configurations, averaged over 50 outer stratified $80/20$ splits. Values are leakage-free outer-test AUCs.}
\label{tab:probe_results}
\end{table}

The raw configuration reaches mean outer-test AUC $0.931$. After linearly removing prompt length from every hidden state dimension, the residualized configuration reaches $0.911$. The prompt-only baseline reaches $0.754$. The gap between the residualized probe and the prompt-only baseline is about $16$ AUC points, placing a concrete floor under the result: a substantial portion of the hidden state's pass/fail signal is not explained by prompt length and survives its removal. The signal is present before a single output token is generated, encoded in the model's internal representation of the prompt alone.

The prompt-only baseline reported above uses logistic regression. To check that this floor does not depend on the choice of a linear model, a random forest, a gradient boosting model, and a small multilayer perceptron were each fit on \texttt{prompt\_tokens} alone under the same nested CV procedure described in Section~\ref{subsec:layer_selection}. Table~\ref{tab:length_baselines} reports their mean outer-test AUC. None exceeds the logistic baseline, and all remain near $0.754$, far below the residualized probe at $0.911$. Prompt length alone does not approach the residualized probe, so the probe's gap over the prompt-length baseline is not an artifact of the baseline being restricted to a linear model.

\begin{table}[h]
\centering
\begin{tabular}{lc}
\toprule
Prompt-only model & Test AUC (mean $\pm$ 95\% CI) \\
\midrule
Logistic regression   & $0.754 \pm 0.014$ \\
Random forest         & $0.735 \pm 0.016$ \\
Gradient boosting     & $0.736 \pm 0.016$ \\
Multilayer perceptron & $0.753 \pm 0.014$ \\
\bottomrule
\end{tabular}
\caption{Mean outer-test AUC for prompt-only models, each fit on \texttt{prompt\_tokens} alone following the nested CV procedure of Section~\ref{subsec:layer_selection}, with model-specific hyperparameter grids tuned in the inner folds. The logistic model is the strongest, and the nonlinear families do not improve on it.}
\label{tab:length_baselines}
\end{table}

\section{Limitation and null result}
\label{sec:limitations}

This section reports one limitation and one null result of the study.

\begin{itemize}
    \item \textbf{Single model.} All experiments use Qwen3-4B-Instruct-2507. Replicating across model families would require substantially more inference compute and time, this is left for future work.

    \item \textbf{Repair geometry is a null result by data starvation.} A companion question asked whether the hidden state shift from a failing attempt to its repair attempt carries a geometric signature of repair success. The data cannot support it. Of the 266 tasks whose first attempt fails, only 20 ever recover within the 5-attempt budget, 14 of them at the first repair, and the conditional pass rate among still-unsolved tasks collapses from 40.1\% at attempt 0 to 5.3\%, 2.0\%, 0.4\% and 0\% at attempts 1 through 4. Prior reports find that iterative self-repair plateaus after the first one or two attempts \citep{arimbur2026howmanytries}, and the decay here is even sharper, with the pass rate collapsing at the very first repair. No contrastive direction can be meaningfully estimated from 14 successful repairs in a 2560-dimensional space, so no claim about repair geometry is made in either direction, and the distinction between ``the signal does not exist'' and ``the data were not sufficient to detect it'' cannot be resolved with the present dataset.
\end{itemize}

\section{Conclusion}
\label{sec:conclusion}

This work examined two related questions about the layer-30 hidden states of Qwen3-4B-Instruct-2507 on LiveCodeBench tasks. 
The first asked whether code correctness is decodable from the prompt-final hidden state, before any token is generated. 
A linear probe answered yes, with performance robust across many random splits and surviving the linear removal of prompt length. 
Prompt length alone, fit with nonlinear models as well as a linear one, stayed near $0.720$ AUC and far below the probe, so its lead is not an artifact of a linear baseline. 
The second asked whether the hidden state shift from a failing attempt to its repair attempt carries a geometric signature of repair success. 
A statistically detectable contrastive direction exists in the raw deltas, and although a conditional residualization against observable repair-context covariates that differ between groups reduces it in both magnitude and stability, the direction remains significant on both tests. 
It is best read as a repair-success signal that the observable repair context does not fully account for.

Both questions in this paper were answered with the help of the same tool, residualization, the operation of regressing each hidden state dimension on a candidate covariate and subtracting the prediction, applied before the downstream test. 
In Section~\ref{sec:probe_results} the probe's signal survived residualizing against prompt length, showing that the hidden states carry a correctness signal not reducible to input length. 
In Section~\ref{sec:repair_geometry} the same operation, this time against three repair-context covariates that differed between groups, removed a substantial part of the geometric direction while leaving a residual that stayed significant on both tests. 
What the raw deltas presented as a single strong direction was in part driven by what differed in the repair prompt itself, yet a genuine repair-success signal remained once that part was subtracted. 
The value of residualization is that it separates these two contributions, and in Section~\ref{sec:repair_geometry} its use was made conditional on whether the covariate to be removed is one that the contrastive construction has not already cancelled, which is what motivated the covariate-distribution check used there.

Several extensions of this work are natural, two of which stand out. The first is to replicate both analyses across other open-weight models, which would help separate findings that are specific to this one model from properties shared across language models. 
The second, more targeted at the repair-direction analysis, is to collect a larger sample of successful multi-step repairs: with more successful transitions at $1 \to 2$, $2 \to 3$, and beyond, the contrastive direction analysis could be repeated step by step rather than only at the first repair. 
Code generation is a particularly tractable domain for this kind of analysis: outputs are objectively verifiable, the failure modes are observable, and each task gives a binary pass/fail signal. More generally, as language models grow more capable their internal computation grows more elaborate, and there is correspondingly more to learn from studying what those internals encode. 
The contribution of this study is as much methodological as empirical. A signal is only as trustworthy as the confound control it survives, and what this paper claims is exactly what survived it.


\bibliographystyle{unsrtnat}
\bibliography{references}

\end{document}
